\subsection{Вспомогательные утверждения: конечномерный случай}

Прежде чем разбирать интегральное уравнение, соберём нужные факты для
конечномерного аналога — обычной системы линейных алгебраических уравнений
того же вида $\vec u - A\vec u = \vec f$.

\medskip
Обозначения: $\vec u = (u_1,\dots,u_m) \in \CC^m$ — вектор-столбец,
$\vec f = (f_1,\dots,f_m) \in \CC^m$, матрица $A \in \CC^{m\times m}$, и
\[
  A^* = \overline{A}^{\,\mathsf{T}}
  \qquad\text{— сопряжённая матрица} .
\]
Скалярное произведение в $\CC^m$:
$(\vec u, \vec v) = \sum_k u_k \overline{v_k}$, так что
$(A\vec u, \vec v) = (\vec u, A^*\vec v)$.

\begin{warn}
В тетради здесь написано «$A : \CC^n \to \CC^m$», $\vec u \in \CC^n$,
$\vec f \in \CC^m$, $A \in \CC^{m\times n}$. Но запись $\vec u - A\vec u$ имеет
смысл только при $n = m$, и все дальнейшие выкладки (определитель, ранг,
$\ker(E-A)$) тоже квадратные. Поэтому здесь всюду $n = m$.
\end{warn}

Рассматриваются три уравнения:
\[
  \vec u - A\vec u = \vec f \qquad (1)
\]
\[
  \vec u - A\vec u = \vec 0 \qquad (1_0)
\]
\[
  \vec v - A^*\vec v = \vec 0 \qquad (1_0^*)
\]
Обозначим
\[
  B = E - A , \qquad B^* = E - A^* .
\]

\subsubsection*{Условие разрешимости}

\begin{proposition}\label{prop:fin-solvable}
Система~(1) разрешима тогда и только тогда, когда $\vec f$ ортогонален всем
решениям сопряжённого однородного уравнения~$(1_0^*)$:
\[
  \vec f \in \big[\ker (E - A^*)\big]^{\perp} .
\]
\end{proposition}

\begin{proof}
Пусть $\vec b_1, \dots, \vec b_m$ — столбцы матрицы $B$. Тогда
\[
  \operatorname{Im} B = \operatorname{span}\{\vec b_1, \dots, \vec b_m\} ,
  \qquad
  \CC^m = \operatorname{span}\{\vec b_k\}
       \ \oplus\ \big(\operatorname{span}\{\vec b_k\}\big)^{\perp} ,
\]
и система~(1) разрешима в точности тогда, когда
$\vec f \in \operatorname{Im} B$.

\medskip
Осталось понять, что такое $(\operatorname{Im} B)^{\perp}$. Вектор $\vec v$
ортогонален образу тогда и только тогда, когда
\[
  (B\vec u, \vec v) = 0 \quad \forall\, \vec u
  \quad\Longleftrightarrow\quad
  (\vec u, B^*\vec v) = 0 \quad \forall\, \vec u
  \quad\Longleftrightarrow\quad
  B^*\vec v = \vec 0 ,
\]
то есть $(\operatorname{Im} B)^{\perp} = \ker B^*$. Значит
$\operatorname{Im} B = (\ker B^*)^{\perp}$, что и требовалось.
\end{proof}

\begin{idea}
Покоординатно последнее равносильно $(\vec b_k, \vec v) = 0$ для всех $k$ —
именно так это и записано в тетради: $\vec v$ ортогонален каждому столбцу
матрицы $B$. Условие $B^*\vec v = 0$ — компактная форма того же набора
равенств.
\end{idea}

\subsubsection*{Число решений однородных уравнений}

\begin{proposition}\label{prop:fin-dims}
Числа линейно независимых решений уравнений $(1_0)$ и $(1_0^*)$ совпадают:
\[
  \dim \ker (E - A) = \dim \ker (E - A^*) .
\]
\end{proposition}

\begin{proof}
Пусть $r = \operatorname{rank} B = \dim \operatorname{Im} B$. Ранг не меняется
ни при транспонировании, ни при комплексном сопряжении, поэтому
$\operatorname{rank} B^* = r$. По теореме о ранге и дефекте
\[
  \dim \ker B = m - r = \dim \ker B^* .
\]
\end{proof}

\subsubsection*{Альтернатива}

\begin{theorem}[альтернатива в конечномерном случае]\label{thm:fin-alternative}
Возможен ровно один из двух случаев.
\begin{enumerate}[label=\arabic*), leftmargin=2em]
  \item $\det(E - A) \ne 0$. Тогда $(1_0)$ имеет только тривиальное решение,
        и система~(1) однозначно разрешима при \emph{любой} правой части
        $\vec f$.
  \item $\det(E - A) = 0$. Тогда $(1_0)$ и $(1_0^*)$ имеют одно и то же
        число $m - r > 0$ линейно независимых решений, система~(1) разрешима
        в точности при $\vec f \perp \ker(E - A^*)$, и её решение определено
        с точностью до прибавления произвольного решения $(1_0)$.
\end{enumerate}
\end{theorem}

\begin{proof}
Случаи взаимно исключают друг друга по определению. В первом случае $B$
обратима — это и есть однозначная разрешимость. Во втором работают
утверждения~\ref{prop:fin-solvable} и~\ref{prop:fin-dims}.

\medskip
Остаётся заметить, что оба однородных уравнения вырождаются одновременно:
\[
  \det B^* = \det\big(\overline{B}^{\,\mathsf{T}}\big)
           = \det \overline{B}
           = \overline{\det B} ,
\]
поэтому $\det B = 0 \iff \det B^* = 0$.
\end{proof}

\begin{idea}
Это и есть \textbf{альтернатива Фредгольма}, только в конечномерном виде, где
она сводится к обычной линейной алгебре. Весь смысл следующего параграфа в
том, чтобы для вырожденного ядра свести интегральное уравнение ровно к такой
системе — и тем самым перенести альтернативу на бесконечномерный случай.
\end{idea}
